\section{Memory map, linker script, and GCC-to-hex flow}

\subsection{Harvard memory map}
The RTL models separate instruction and data memories. Software is linked to
match the address split so code and writable data never overlap.

\begin{center}
\begin{tabularx}{\linewidth}{@{}l l l l X@{}}
\toprule
Region & Base address & Size & Access & Purpose \\
\midrule
IMEM & \texttt{0x0000\_0000} & 8 KiB & read/execute & \texttt{.text}, \texttt{.init}, \texttt{.rodata} \\
DMEM & \texttt{0x2000\_0000} & 8 KiB & read/write & \texttt{.data}, \texttt{.sdata}, \texttt{.bss}, stack, signature window \\
\bottomrule
\end{tabularx}
\end{center}

Tests use a fixed signature scratch region at \texttt{0x2000\_0100}: programs
write deterministic 32-bit words there, and the testbench compares them against
golden \texttt{.sig} files.

\subsection{Why these addresses?}
The documentation emphasizes a few motivations:
\begin{itemize}[leftmargin=*]
  \item The PC resets to zero, so placing IMEM at \texttt{0x0000\_0000} makes reset-to-first-instruction trivial.
  \item The DMEM base at \texttt{0x2000\_0000} makes the Harvard split obvious from the address, and leaves room to grow toward a small SoC/MMIO map.
  \item The linker fixes a deterministic stack top and a deterministic signature window for self-checking tests.
\end{itemize}

\subsection{A note on address decoding in the simple memories}
The VHDL \texttt{imem\_dp} and \texttt{dmem} models intentionally do not
implement full high-bit address decode; they mask down to the configured depth
when indexing their arrays. The
software \emph{still} uses full logical addresses (IMEM at \texttt{0x0000\_xxxx},
DMEM at \texttt{0x2000\_xxxx}) so that the linker keeps sections separated.
When integrating into a larger system, you can wrap proper decode logic around
the memories without changing the compiler/linker view.

\subsection{Data-side reads from the text region (\texttt{.rodata})}
Because \texttt{.rodata} constants are placed in IMEM alongside \texttt{.text},
loads that target addresses below \texttt{DMEM\_BASE\_ADDR}
(\texttt{0x2000\_0000}) must read from IMEM, not DMEM. The top-level wrapper
\texttt{top.vhd} instantiates a true dual-port instruction memory
(\texttt{imem\_dp}): port A serves instruction fetch and port B serves data-side
reads. A registered mux-select (\texttt{use\_imem\_w}) aligns the IMEM/DMEM
read-data mux with the synchronous BRAM output timing.

\subsection{Linker script (\texttt{link.ld})}
The provided linker script enforces the memory map:
\begin{itemize}[leftmargin=*]
  \item \texttt{.text} and \texttt{.rodata} go to IMEM at \texttt{0x0000\_0000} (the PC resets to zero).
  \item \texttt{.data} (and small-data \texttt{.sdata}) live in DMEM at runtime but are placed \emph{in the ELF load image} so the build can initialize DMEM before releasing reset.
  \item \texttt{.bss} (and \texttt{.sbss}) are zero-initialized at runtime and do not need hex initialization data.
  \item \texttt{\_\_stack\_top} is set to the top of DMEM so freestanding startup code can initialize \texttt{sp}.
  \item \texttt{.note} sections are discarded, and the link stops with an error if \texttt{.text} does not start at \texttt{0x0000\_0000}. Some toolchains, for example the RISC-V GCC that comes with AMD Vitis, add a build-ID note section by default; placed in front of \texttt{.text}, it would move the code away from the address where the IMEM image starts.
\end{itemize}

\paragraph{Key snippet}
\begin{lstlisting}[style=minrv]
MEMORY {
  IMEM (rx) : ORIGIN = 0x00000000, LENGTH = 8K
  DMEM (rw) : ORIGIN = 0x20000000, LENGTH = 8K
}
\end{lstlisting}

\paragraph{Spike variant (\texttt{link\_spike.ld})}
The golden signatures are generated on the Spike ISA simulator
(Section~\ref{sec:spike-golden}). Spike reserves the low addresses for its debug
module and boot ROM, so \texttt{link\_spike.ld} places the code at
\texttt{0x8000\_0000} and keeps DMEM, the stack, and the signature window at
their MinRV32 addresses. \texttt{.data} is placed at its run address, so Spike's
ELF loader preloads DMEM as \texttt{<name>.data.hex} does on the core.
\texttt{build\_prog.sh} selects it through the \texttt{LINKER\_SCRIPT}
environment variable. Section~\ref{sec:spike-addr} explains which signature
values this affects.

\subsection{GCC $\rightarrow$ ELF $\rightarrow$ hex images}
The project's workflow emphasizes real software running on the RTL:
\begin{enumerate}[leftmargin=*]
  \item Compile/link a freestanding RV32I program with a bare-metal RISC-V GCC toolchain.
  \item Extract \texttt{.text/.rodata} and \texttt{.data/.sdata} into raw binaries via \texttt{objcopy}.
  \item Convert those binaries into newline-delimited 32-bit hex words (\texttt{XXXXXXXX} per line) for IMEM/DMEM initialization.
\end{enumerate}

The helper script \texttt{scripts/build\_prog.sh} automates this and produces:
\texttt{<base>.elf}, \texttt{<base>.text.hex}, \texttt{<base>.data.hex} (and, for
C sources, a compiler-generated assembly listing \texttt{<base>.S}).

\paragraph{Practical guard: do the images fit?}
The educational memories are small (default 8\,KiB IMEM and 8\,KiB DMEM). The build script
prints a size summary for the generated \texttt{.text} and \texttt{.data} images
and emits warnings if either exceeds the configured IMEM/DMEM capacity. If you
change the VHDL memory sizes, set \texttt{IMEM\_BYTES} and \texttt{DMEM\_BYTES}
in the environment so the warnings match your configuration.

\subsection{Future work: one memory map for the core and Spike}
\label{sec:future-memmap}
The two linker scripts exist only because the core's IMEM sits at
\texttt{0x0}, where Spike has its debug module. Spike does not accept memory
there. Its other devices (boot ROM at \texttt{0x1000}, CLINT at
\texttt{0x0200\_0000}, PLIC at \texttt{0x0C00\_0000}, UART at
\texttt{0x1000\_0000}) do not overlap DMEM at \texttt{0x2000\_0000}. If IMEM
and the reset PC move to \texttt{0x8000\_0000} (the RAM base also used by
riscv-tests, riscv-arch-test, QEMU \texttt{virt}, and Sail), one linker script
serves both. DMEM, the stack, and the signature window do not change. The
current layout does not affect the golden signatures as long as tests follow
the code-address rule (Section~\ref{sec:spike-addr}), so this change is not
needed now. Required changes:
\begin{enumerate}[leftmargin=*]
  \item \textbf{Reset PC.} Add a reset-value generic to \texttt{pipe\_reg}
  (default zero) and a \texttt{RESET\_PC} generic (default
  \texttt{x"80000000"}) passed down to \texttt{pc\_reg} in
  \texttt{fetch\_stage}.
  \item \textbf{Load-path decode.} \texttt{top.vhd} and
  \texttt{sim/top\_trace.vhd} send a load to IMEM port~B when
  \texttt{addr < DMEM\_BASE\_ADDR}; test address bit~31 instead (set for
  IMEM, clear for DMEM), which also needs less logic.
  \item \textbf{Memories.} No change: \texttt{imem\_dp} and \texttt{dmem} index
  only the low address bits.
  \item \textbf{Linker script.} Set \texttt{IMEM ORIGIN = 0x80000000} in
  \texttt{link.ld} and remove \texttt{AT>IMEM} (nothing copies \texttt{.data}
  from its load address). Delete \texttt{link\_spike.ld} and the
  \texttt{LINKER\_SCRIPT} override in \texttt{gen\_golden.sh}.
  \item \textbf{Documentation.} Update the address tables, remove the
  code-address rule, and update examples that show PC values.
\end{enumerate}
Raw code addresses would then give the same signature on both, and the existing
\texttt{.sig} files stay valid. The PC would no longer reset to zero. The
differences in Section~\ref{sec:spike-other} remain.
